\chapter{CAD-Aufbau}
\label{ch:cad}

\section{Konzeptentwicklung}
\label{sec:concept}

\section{Flügelsteuermechanismus}
\label{sec:fin_control}
Nach dem Entwurf der gesamten Raketengeometrie und der Festlegung der Form/Position
der Flügel, wurde die Entwicklung eines Steuermechanismus für die Flügel begonnen.
Nach ein Paar gescheiterten Versuchen, wurde ein Mechanismus entwickelt, bei dem
sowohl die Motoren aufgrund von Platzmangel in das Gehäuse passten, als diese auch
effizient das Drehmoment auf die Flügel übertragen konnten. 
\\
Aufgrund von Platzmangel im Heckteil der Rakete, konnten die die Servomotoren nicht
direkt an die Flügel angebracht werden, sondern mussten über ein System mit 2 ineinander
zeigenden Hebelarmen und einem Gelenkarm, der die zwei Hebelarme miteinander verbindet.
Das Drehmoment der Servomotoren wird dabei 1 zu 1 übertragen, da beide Hebelarme die
gleiche Länge einer Drehachse auf die andere haben. Abbildung \ref{fig:fin_servo} zeigt
den Flügelsteuermechanismus von ``ALBERT''.
\begin{figure}[H]
    \centering
        \includegraphics[width=0.4\textwidth]{fin_servo.png}
        \caption{Flügelsteuermechanismus}
        \label{fig:fin_servo}
\end{figure}
Um sicherzustellen, dass die Flügel stabil am Gehäuse befestigt sind, wurden diese mit
2 Kugellagern gelagert, die in das Gehäuse eingebaut wurden. 

\section{Innenstruktur und Materialwahl}
\label{sec:internal_structure}
Die Innenstruktur der Rakete wird auf 4 Carbonstäbe aufgebaut, die vom Heck bis zum
Fallschirmfach verlaufen. Diese Stäbe sorgen für die notwendige Steifigkeit der Rakete
und helfen dabei, die eizelnen Komponenten in nötiger Position zu halten. Die Stäbe
mussten allerdings im Bereich des Flightcomputers unterbrochen werden, um Platz für
diesen zu schaffen.
\\
\\
Die Wahl des Materials für die Rakete fiel auf 3D-gedrucktes PLA, da dieses Material
leicht zu verarbeiten ist, eine ausreichende Festigkeit für die Anforderungen der Rakete
bietet und kostengünstig ist. Allerdings um die Festigkeit der Flügel und dessen
Steuermechanismus zu erhöhen, wurde das Greentec Pro Filament von Extrudr verwendet,
da dieses Material genauso einfach wie PLA zu verarbeiten ist, aber eine deutlich höhere
Festigkeit aufweist.
\\
\\
Für das Außengehäuse der Rakete wurde ein Kartorohr mit einem Außendurchmesser von 84 mm
und einer Wandstärke von 2mm verwendet. Dieses ist eine kostengünstige und leicht verfügbare
Option, die alle Anforderungen an die Festigkeit und Steifigkeit der Rakete ausreichend erfüllt.

\section{Implementierung und Fertigung}
\label{sec:implementation}
Die Fertigung der Rakete erfolgte hauptsächlich durch 3D-Druck, wobei die meisten Komponenten
mit dem Bambulab A1 mini gedruckt wurden. Die einzige Ausnahme war die Raketenspitze,
die aufgrund ihrer Größe auf einem größeren Drucker gedruckt werden musste. 

\section{Herausforderungen und Lösungen}
\label{sec:challenges}

\section{Fazit}
\label{sec:cad_fazit}

%=====================================================%

\chapter{Simulation der Aerodynamik}
\label{ch:aero_simulation}

\section{Simulation des Raketenkörpers in Ansys Fluent}
\label{sec:used_software_sim}
Um alle nötigen Kraft-/Momentkoeffizienten für das Steuersystem zu berechnen, wurde der ganze Raketenkörper in Ansys Fluent mit Verschiedenen Anstellwinkeln der Flügel simuliert.

\subsection{Definition der Referenzwerte und der simulierten/berechneten Werte}
\label{subsec:reference_sim_values}
Um die richtige Berechnung von den Kraft-/Momentkoeffizienten und dessen Ableitungen müssen in Fluent bestimmte Referenzwerte gesetzt werden.
\paragraph{Definition der Referenzwerte für alle Simulationsdurchgänge}
\begin{itemize}
    \item $A = \frac{\pi \cdot D^2}{4}$: Referenzfläche (Frontfläche der Rakete) mit $D$ als Durchmesser des Raketenkörpers
    \item $l = D$: Referenzlänge (Durchmesser des Raketenkörpers)
\end{itemize}
\paragraph{Definition der Winkel des Raketenkörpers und der Flügel}
\begin{itemize}
    \item $\alpha$: Pitch Winkel des Raketenkörpers
    \item $\beta$: Yaw Winkel des Raketenkörpers
    \item $\delta_i$: Anstellwinkel von Flügel $i$
    \item $\eta$: Anstellwinkel vom für Pitch verantwortliches Flügelpaar
    \item $\zeta$: Anstellwinkeln vom für Yaw verantwortliches Flügelpaar
\end{itemize}
\paragraph{Aus der Simulation bezogene Kräfte/Beiwerte}
\begin{itemize}
    \item $F_A$: Axialkraft
    \item $F_Y$: Seitenkraft
    \item $F_Z$: Normalkraft
    \item $C_l$: Rollmomentbeiwert
    \item $C_m$: Pitchmomentbeiwert
    \item $C_n$: Yawmomentbeiwert
\end{itemize}

\subsection{Simulationsdurchgänge}
\label{subsec:sim_cases}
Um genug Messpunkte der verschiedenen Werte für genaue Ableitungen zu erreichen und um für alle simulierten Werte die besten Bedingungen zu schaffen, wo diese am genausten von Fluent berechent werden können, wurden 11 Verschieden Fälle in der Simulation betrachtet. Durch diese 11 Durchgänge wurden sowohl alle möglichen Fälle bezüglich Pitch- und Yaw-Winkel abgedeckt, als auch Veränderungen im Anstellwinkel der Flügel. Da der Raketenkörper symmetrisch ist, sind die Beiwerte für Pitch und Yaw Stellungen identisch. Daher können die in der Simulation errechneten Werte für Pitch Änderung auch für Yaw-bezogene Beiwerte verwendet werden. Diese Werte für $\alpha, \beta, \delta_i$ bzw. $\eta$ und $\zeta$ wurden bei jedem Durchgang angenommen:
\begin{itemize}
    \item Fall 1: $\alpha = \ang{0}, \beta = \ang{0}, \eta = \ang{0}, \zeta = \ang{0}$
    \item Fall 2: $\alpha = \ang{0}, \beta = \ang{0}, \eta = \ang{2}, \zeta = \ang{0}$
    \item Fall 3: $\alpha = \ang{0}, \beta = \ang{0}, \eta = \ang{5}, \zeta = \ang{0}$
    \item Fall 4: $\alpha = \ang{2}, \beta = \ang{0}, \eta = \ang{0}, \zeta = \ang{0}$
    \item Fall 5: $\alpha = \ang{2}, \beta = \ang{0}, \eta = \ang{2}, \zeta = \ang{0}$
    \item Fall 6: $\alpha = \ang{2}, \beta = \ang{0}, \eta = \ang{5}, \zeta = \ang{0}$
    \item Fall 7: $\alpha = \ang{5}, \beta = \ang{0}, \eta = \ang{0}, \zeta = \ang{0}$
    \item Fall 8: $\alpha = \ang{5}, \beta = \ang{0}, \eta = \ang{2}, \zeta = \ang{0}$
    \item Fall 9: $\alpha = \ang{5}, \beta = \ang{0}, \eta = \ang{5}, \zeta = \ang{0}$
    \item Fall 10: $\alpha = \ang{5}, \beta = \ang{0}, \delta_1 = \ang{2}, \delta_2 = \ang{2}, \delta_3 = \ang{2}, \delta_4 = \ang{2}$
    \item Fall 11: $\alpha = \ang{0}, \beta = \ang{0}, \delta_1 = \ang{5}, \delta_2 = \ang{5}, \delta_3 = \ang{5}, \delta_4 = \ang{5}$
\end{itemize}

\section{Auswertung der Simulationsergebnisse}
\label{sec:simulation_results}
Nach dem Simulieren von allen 11 Fällen, wurden diese entsprechend ausgewertet und alle für die Regelung nötigen Beiwerte wurden errechnet. Im ersten Schritt wurden alle aus den Simulationsergebnissen bezogene Kräfte in entsprechende Beiwerte umgerechnet.

\subsection*{Berechnung der Kraftbeiwerte}
Da es in der Student-Version von Ansys Fluent nicht möglich ist bei einem Kraft-Report dessen Beiwert zu berechnen, müssen die Axial-, Normal- und Seitkräfte mit der Hand auf die jeweiligen Beiwerte umgerechnet werden. Kraftbeiwerte werden grundsätzlich nach folgenden Prinzip berechnet:
\begin{equation}
    C = \frac{F}{q \cdot A_{ref}}
\end{equation}
Dabei ist $C$ der resultierende Beiwert, $F$ die simulierte Kraft, $A_{ref}$ die Referenzfläche und $q$ der dynamische Druck der wie folgt definiert ist:
\begin{equation}
    q = \frac{1}{2}\rho v^2
\end{equation}

\subsection*{Simulationsergebnisse und Plots} 
\label{subsec:sim_results_plots}

\subsection*{Berechnungen der ``Modellwerte'' und Ableitungen}
\label{subsec:derivatives_calculation}
Als ersten Schritt wurden die Widerstandbeiwerte berechnet. Diese resultieren aus den folgenden Ableitungen:
\begin{equation}
    C_A = C_{A0} + \frac{\Delta C_A}{\Delta \alpha} + \frac{\Delta C_A}{\Delta \beta}
\end{equation}
Hier wird der Axialkraftbeiwert $C_A$ aus dem Axialkraftbeiwert ohne jeglichen Anstellwinkeln $C_{A0}$, der Ableitung des Axialkraftbeiwertes nach Pitch Winkel $\frac{\Delta C_A}{\Delta \alpha}$ und der Ableitung nach Yaw Winkel $\frac{\Delta C_A}{\Delta \beta}$.
\begin{equation}
    C_Z = \frac{\Delta C_Z}{\Delta \alpha} + \frac{\Delta C_Z}{\Delta \eta}
\end{equation}
Hierbei resultiert der Normalkraftbeiwert $C_Z$ aus dessen Ableitung nach Pitch Winkel $\frac{\Delta C_Z}{\Delta \alpha}$ und nach Anstellwinkel des für Pitch verantwortliches Flügelpaar $\frac{\Delta C_Z}{\Delta \eta}$.
\begin{equation}
    C_Y = \frac{\Delta C_Y}{\Delta \beta} + \frac{\Delta C_Y}{\Delta \zeta}
\end{equation}
Anschließend ergibt sich der Seitenkraftbeiwert $C_Y$ aus dessen Ableitung nach Yaw Winkel $\frac{\Delta C_Y}{\Delta \beta}$ und nach dem Anstellwinkel des für Yaw verantwortliches Flügelpaar $\frac{\Delta C_Y}{\Delta \zeta}$.
\\


\subsubsection*{Visualisierung der Strömung um den Raketenkörper}
\label{subsec:cfd_flow_visualization}

\section{Messung eines Raketenflügels im Windkanal}
\label{sec:wind_tunnel_test}

\section{Vergleich mit Messdaten}
\label{sec:comparison_wind_tunnel_cfd}

\section{Herausforderungen und Lösungen}
\label{sec:challenges_simulation}

\section{Fazit}
\label{sec:simulation_fazit}

